% from ega4-2-fr.tex
\oldpage[IV-2]{62}

\begin{proof}
En effet, pour toute extension $K$ de $k$, $f_{(K)}:X_{(K)}\to Y_{(K)}$ est
surjectif (I, 3.5.2), et $X_{(K)}$ étant irréductible (resp. connexe) par
hypothèse, il en est de même de $Y_{(K)}$.
\end{proof}

\begin{definition}[4.5.5]
\label{IV.2.4.5.5-fr}
On dit qu'un morphisme $f:X\to Y$ de préschémas est \emph{irréductible}
(resp. \emph{connexe}), si, pour tout $y\in Y$, le $k(y)$-préschéma
$f^{-1}(y)$ est géométriquement irréductible (resp. géométriquement connexe).
\end{definition}

\begin{proposition}[4.5.6]
\label{IV.2.4.5.6-fr}
\begin{enumerate}
\item[(i)] Soient $X$ un $k$-préschéma, $K$ une extension de $k$. Alors le
nombre géométrique de composantes irréductibles (resp. connexes) de $X_{(K)}$
relativement à $K$ est égal au nombre géométrique de composantes irréductibles
(resp. connexes) de $X$ relativement à $k$. En particulier, pour que le
$K$-préschéma $X_{(K)}$ soit géométriquement irréductible (resp.
géométriquement connexe), il faut et il suffit que le $k$-préschéma $X$ le
soit.
\item[(ii)] Soient $f:X\to Y$, $g:Y'\to Y$ deux morphismes, et posons
$X'=X_{(Y')}=X\times_Y Y'$, $f'=f_{(Y')}:X'\to Y'$. Si $f$ est irréductible
(resp. connexe), il en est de même de $f'$. La réciproque est vraie si $g$ est
surjectif.
\end{enumerate}
\end{proposition}

\begin{proof}
Pour (i), si $\Omega$ est une extension algébriquement close de $K$, on a
$X\otimes_k\Omega=X_{(K)}\otimes_K\Omega$, d'où l'assertion.

Pour (ii), pour tout $y'\in Y'$, si l'on pose $y=g(y')$, on a
$f'^{-1}(y')=f^{-1}(y)\otimes_{k(y)}k(y')$ (I, 3.6.4), donc les deux assertions
résultent de (i).
\end{proof}

\begin{proposition}[4.5.7]
\label{IV.2.4.5.7-fr}
Soient $f:X\to Y$ un morphisme surjectif irréductible (resp. connexe),
$g:Y'\to Y$ un morphisme, et posons $X'=X\times_Y Y'$. Si en outre $Y'$ est
irréductible (resp. connexe), et $f$ universellement ouvert (resp. plat et
quasi-compact, ou universellement ouvert, ou universellement fermé), alors
$X'$ est irréductible (resp. connexe).
\end{proposition}

\begin{proof}
Les propriétés de $f$ que l'on considère étant toutes stables par changement
de base, on peut se borner au cas où $Y'=Y$~; il suffit alors d'appliquer
(4.4.2) (compte tenu de (2.3.12)).
\end{proof}

\begin{corollary}[4.5.8]
\label{IV.2.4.5.8-fr}
Soient $X$, $Y$ deux $k$-préschémas.
\begin{enumerate}
\item[(i)] Si $X$ est géométriquement irréductible (resp. géométriquement
connexe) et $Y$ irréductible (resp. connexe), alors $X\times_k Y$ est
irréductible (resp. connexe).
\item[(ii)] Si $X$ et $Y$ sont géométriquement irréductibles (resp.
géométriquement connexes), il en est de même de $X\times_k Y$.
\end{enumerate}
\end{corollary}

\begin{proof}
Pour (i), le morphisme structural $X\to\operatorname{Spec}(k)$ est surjectif et
universellement ouvert (2.4.9) et en outre irréductible (resp. connexe)~; il
suffit donc d'appliquer (4.5.7).

Pour (ii), si $\Omega$ est une extension algébriquement close de $k$, on a
$(X\times_k Y)_{(\Omega)}=X_{(\Omega)}\times_\Omega Y_{(\Omega)}$
(I, 3.3.10)~; par hypothèse $X_{(\Omega)}$ et $Y_{(\Omega)}$ sont des
$\Omega$-préschémas géométriquement irréductibles (resp. géométriquement
connexes) (4.5.6) et il suffit donc d'appliquer (i).
\end{proof}

\begin{proposition}[4.5.9]
\label{IV.2.4.5.9-fr}
Soit $X$ un $k$-préschéma. Les conditions suivantes sont équivalentes~:
\begin{enumerate}
\item[$a)$] $X$ est géométriquement irréductible (autrement dit, pour toute
extension $K$ de $k$, $X\otimes_k K$ est irréductible).
\item[$b)$] Pour toute extension finie séparable $K$ de $k$, $X\otimes_k K$
est irréductible.
\item[$c)$] $X$ est irréductible, et si $x$ est son point générique, $k(x)$
est une extension primaire de $k$.
\end{enumerate}
\end{proposition}

\begin{proof}
Il est clair que $a)$ entraîne $b)$. Pour voir que $b)$ entraîne $c)$,
considérons une extension finie séparable $K$ de $k$~; si
$p:X\otimes_k K\to X$ est la projection canonique, les

